\documentclass[12pt,a4paper]{article}
% ============================================================
%  Structural Persistence Series — Shared Preamble
% ============================================================
\usepackage{fontspec}
\setmainfont{Hiragino Mincho ProN}
\setsansfont{Hiragino Sans}
\setmonofont{Menlo}

\usepackage{iftex}
\ifXeTeX
  \XeTeXlinebreaklocale "ja"
  \XeTeXlinebreakskip=0pt plus 1pt minus 0.1pt
\fi

\usepackage[margin=22mm]{geometry}
\usepackage{amsmath,amssymb,amsthm}
\usepackage{xcolor}
\usepackage{graphicx}
\usepackage{hyperref}
\usepackage{booktabs}
\usepackage{fancyhdr}
% enumitem not available in this TeX installation

% --- Colours ------------------------------------------------
\definecolor{PaperAccent}{HTML}{1F4E79}
\definecolor{PaperGray}{HTML}{51606F}

\hypersetup{
  colorlinks=true,
  linkcolor=PaperAccent,
  urlcolor=PaperAccent,
  citecolor=PaperAccent,
  pdflang=ja
}
\urlstyle{same}

% --- Typography ---------------------------------------------
\setlength{\parindent}{1.5em}
\setlength{\parskip}{0pt}
\setlength{\emergencystretch}{3em}
\linespread{1.08}
\setlength{\abovedisplayskip}{8pt plus 2pt minus 4pt}
\setlength{\belowdisplayskip}{8pt plus 2pt minus 4pt}
\setlength{\abovedisplayshortskip}{6pt plus 2pt minus 2pt}
\setlength{\belowdisplayshortskip}{6pt plus 2pt minus 2pt}

% --- Section label names ------------------------------------
\renewcommand{\abstractname}{要旨}

% --- Theorem-like environments ------------------------------
\theoremstyle{plain}
\newtheorem{proposition}{命題}[section]
\newtheorem{lemma}[proposition]{補題}
\newtheorem{corollary}[proposition]{系}

\theoremstyle{definition}
\newtheorem{definition}{定義}[section]
\newtheorem{assumption}{条件}[section]

\theoremstyle{remark}
\newtheorem*{remark}{注}

% --- Running header -----------------------------------------
\pagestyle{fancy}
\fancyhf{}
\renewcommand{\headrulewidth}{0.4pt}
\fancyhead[L]{\small\textcolor{PaperGray}{\nouppercase{\leftmark}}}
\fancyhead[R]{\small\textcolor{PaperGray}{\thepage}}
\fancypagestyle{plain}{%
  \fancyhf{}
  \renewcommand{\headrulewidth}{0pt}
  \fancyfoot[C]{\small\textcolor{PaperGray}{\thepage}}
}

% --- Title block macro --------------------------------------
\newcommand{\PaperTitleBlock}[3]{%
  % #1 = title, #2 = subtitle, #3 = date
  \thispagestyle{plain}%
  \begin{center}
  {\LARGE\bfseries #1\par}
  \vspace{0.4em}
  {\normalsize #2\par}
  \vspace{1.0em}
  {\large Akihito Sunagawa\par}
  \vspace{0.3em}
  {\normalsize #3\par}
  \end{center}
  \vspace{1.6em}
}

% Legacy 2-arg version (backward compat)
\newcommand{\PaperTitleBlockSimple}[2]{%
  \PaperTitleBlock{#1}{}{#2}
}


\begin{document}

\PaperTitleBlock
  {構造持続の条件つき導出}
  {— 条件つき導出と弱依存下の安定性 —}
  {2026年4月12日}

\begin{abstract}
前稿で示した構造持続の最小形式 \(S = Me^{-L}\) の背後には、状態集合の残存量に関する指数式 \(m(V^{(n)}) = m(V^{(0)}) e^{-L}\) がある。この式は、状態集合の縮小率を対数比で定義する限り恒等式として成り立ち、損失どうしの独立性を要しない。本稿の目的は、この恒等式を繰り返すことではなく、どこまでが定義の帰結であり、どこからが追加条件に依存するのかを明示することにある。

そのために、まず A1--A2 のもとで恒等式が成り立つことを確認する。ついで、段階損失の生成過程を確率的にモデル化する際の独立性条件 A3 を導入し、その緩和のもとでどこまで安定性が保たれるかを述べる。最後に、Lean による形式検証の範囲を記す。本稿の役割は、最小形式の背後にある数学的骨格を、必要最小限のかたちで与えることにある。
\end{abstract}

\section{はじめに}\label{ux306fux3058ux3081ux306b}

最小形式の中心式 \[
S = Me^{-L}
\]

は、構造持続を記述するうえで、できるだけ仮定の少ない形を与えると同時に、広い応用可能性を持つ。しかし、その薄さゆえに、この指数型がどこまで定義の帰結であり、どこから追加条件に依存するのかは見えにくい。

本稿の目的は、この点を明示することである。ここで行うのは、最小形式そのものを拡張することではない。むしろ、その背後にある条件つきの導出を、必要最小限のかたちで切り出すことである。

以下では、まず A1--A2 という二つの条件のもとで指数式が恒等式として成り立つことを確認する。ついで、段階損失の生成過程を確率的にモデル化する際の独立性条件 A3 を導入し、その厳密版および弱依存への緩和版のもとでの安定性を述べる。最後に、既存の指数表現と弱依存境界、および A2 の特徴づけ定理（Paper 1 §3 の対数比の一意性定理）が Lean によりどこまで形式検証されているかを記す。

\section{三つの条件}\label{ux4e09ux3064ux306eux6761ux4ef6}

本稿で用いる条件は、次の三つである。

\begin{assumption}[A1]
各段階で加わる制約は、構造を維持できる状態の集合を広げず、縮めるか、そのままに保つ。
\end{assumption}

\begin{assumption}[A2]
各段階の損失は、その段階で生じた縮小率の対数比として測られる。すなわち
$l_i = -\ln(m(V^{(i)})/m(V^{(i-1)}))$。

この関数形は任意の選択ではなく、Paper 1 §3 の対数比の一意性定理（公理系 B1–B4: 比率依存性、正規化、加法性、連続性）の帰結として一意に強制される。B3（加法性）は集合論的には連鎖加法性
$D(A,C) = D(A,B) + D(B,C)$
に、独立分解形では示量性（独立サブシステムでの損失加法性）に対応し、Hartley $(1928)$ / Shannon $(1948)$ による情報量の加法性公理化と同じ骨格を持つ動機づけを備える。
\end{assumption}

\begin{assumption}[A3]
制約の生成過程を確率的にモデル化するとき、各段階の縮小率 $R_i := m(V^{(i)})/m(V^{(i-1)})$ を確率変数と見なし、対応する段階損失 $l_i = -\ln R_i$ について、各 $l_i$ はそれ以前の損失列 ($l_1$, ..., $l_{i-1}$) から独立である。
\end{assumption}

A1 は、制約が縮小方向に働くという条件である。Paper 1 §2 は、この A1 に対する自然な十分条件として、各段階が制約追加モデル \[
V^{(i)} = V^{(i-1)} ∩ C_{i}
\]

で表されるなら縮小列が命題として従うことを示している。本稿では、修復・学習・外部支援・ロールバックのような再拡大モードも論理的には排除しないため、A1 自体はなお一般形の適用条件として保持する。A2 は、損失の測り方を定める条件であるが、Paper 1 §3 の対数比の一意性定理により、これは任意の定義ではなく、構造損失測度に対する自然な公理系 B1--B4 の帰結として一意に強制される。したがって本稿における A1--A2 は、A1 の一つと、Paper 1 §3 の公理系から導かれる A2 の一つからなる二条件であり、両者はそれぞれ独立の内容を持つ。A1--A2 のもとで、3節の指数式は恒等式として成り立つ。A3 はこの恒等式そのものには不要であり、4節で弱依存下の安定性を論じるために導入する。

ここで重要なのは、A2 の背後にある B3（加法性）と A3（段階損失の確率的独立性）が異なる層の独立性条件であることである。B3 は \(f\) の関数方程式であり、その背後の集合論的解釈は集合列の入れ子構造 \(A \supseteq B \supseteq C\) に対する連鎖加法性、あるいは状態空間が二つの独立なサブシステムの積に分解されたときの示量性である。これに対し A3 は、段階損失の生成過程を確率的にモデル化する際に現れる統計的独立性である。両者は直交しており、混同してはならない。

A3 は、実系一般でそのまま成立する経験法則としてではなく、近似的成立を論じるための基準として導入される。厳密独立が失われたときに理論が直ちに無意味になるのではなく、主張の強さが等式から境界へ、さらには個別ドメインでのより弱い実証的主張へと移ることを示すのが、本稿の目的である（4節参照）。

したがって、本稿で扱う主張は二つに分かれる。第一に、対数比定義から従う恒等式としての指数表現である。第二に、段階損失の生成過程を確率的にモデル化したとき、その独立性または弱依存のもとで指数型の安定性がどこまで保たれるか、という条件つき結果である。

\section{恒等式としての指数表現}\label{ux6052ux7b49ux5f0fux3068ux3057ux3066ux306eux6307ux6570ux8868ux73fe}

構造を維持できる状態の残存量は、累積損失に対して指数型で表される。すなわち、初期残存量を \(m(V^{(0)})\)、累積損失を \(L\) とすると、 \[
m(V^{(n)}) = m(V^{(0)}) e^{-L}
\]

が恒等式として成り立つ。これは前稿の命題1と同一であり、A3（独立性）を要しない。恒等式そのものは、A2 の関数形の特徴づけ（Paper 1 §3 の対数比の一意性定理）と累積損失の望遠鏡積から従う。A1 は各段階損失の非負性（\(l_i \ge 0\)）を保証する。集合列の縮小率を対数比で定義し、その和を累積損失とする限り、望遠鏡積によって自動的に得られる。

したがって、指数型は経験的に当てはめた関数形ではなく、公理系から二段階で導かれる帰結である。第一段階として、A2 の関数形（対数比）は Paper 1 §3 の対数比の一意性定理によって公理的に特徴づけられる。第二段階として、この対数比損失のもとで累積損失の望遠鏡積から指数表現が恒等的に従う。最小形式における \(S = M e^{-L}\) は、この残存量の指数表現に、有効維持資源 \(M\) を別に導入した拡張的定義である。

本稿の新規性は、この恒等式そのものにではなく、次節で述べる弱依存下の安定性にある。

\section{近似条件のもとでの安定性}\label{ux8fd1ux4f3cux6761ux4ef6ux306eux3082ux3068ux3067ux306eux5b89ux5b9aux6027}

現実の系では、A3 の意味での厳密な独立が常に成立するとは限らない。以下では、A3 を弱依存の条件へ緩和した場合を考える。

3節の恒等式は、実現された集合列に対しては常に成り立つ。ここで問うのは、段階損失の生成過程に弱い依存があるとき、累積損失がどの範囲に収まるかである。

段階損失 \(l_i\) が確率変数として与えられるとき、各段階損失が独立に加算される参照モデルにおける累積損失を \(L_{\text{ref}}\) とする。依存を含む実際の生成過程のもとで生じる残存可能性を \(P\)、対応する実現累積損失を \(\tilde{L} := -\ln P\) とおく。ここで、依存の効果が参照モデルからの相対誤差として \(\rho\)（\(0 \le \rho < 1\)）で抑えられているとする。すなわち、 \[
(1-\rho)L_{\text{ref}} \le \tilde{L} \le (1+\rho)L_{\text{ref}}
\]

が成り立つならば、同値に \[
e^{-L_{\text{ref}}(1+\rho)} \le P \le e^{-L_{\text{ref}}(1-\rho)}
\]

である。\(\rho = 0\) のとき \(P = e^{-L_{\text{ref}}}\) に戻る。すなわち、実現累積損失 \(\tilde{L}\) が参照量 \(L_{\text{ref}}\) から相対誤差 \(\rho\) の範囲に収まる限り、指数型の記述は保たれる。

ここで \(\rho\) は、段階損失間の依存が累積損失に及ぼす影響の大きさを制御する定数である。\(\rho\) が 0 に近いほど独立ケースに近く、本稿ではそのような場合を弱依存と呼ぶ。

したがって、厳密な独立が成立しなくても、段階損失間の依存が累積損失の相対誤差として \(\rho\) で抑えられる限り、結論は指数境界として安定に保たれる。この点は重要である。実際の系への適用では、厳密な独立性そのものよりも、その近似的成立が意味を持つ場合が多いからである。

なお、本節の結果は弱依存の一般理論を与えるものではない。与えられるのは、依存の効果が累積損失の相対誤差として要約できる場合に対する安定性結果である。個別の系において、具体的な依存構造からこの相対誤差制御を導くことは、応用側の課題として残される。

\section{形式検証}\label{ux5f62ux5f0fux691cux8a3c}

A1--A3 を含む公理環境のもとでの指数表現の導出、および A3 の弱依存緩和に対する指数境界は、Lean により形式検証されている。また、Paper 1 §3 の対数比の一意性定理（B1--B4 からの A2 の特徴づけ）についても、独立の Lean ファイルとして形式検証されている。3節の望遠鏡積による恒等式（A1--A2 のみ）は、これら二つの独立した検証を組み合わせることで得られる。

ここで形式検証の対象となっているのは、明示した定義と仮定のもとでの数学的導出そのものである。したがって、保証されるのは条件つき結果の論理的一貫性であり、各条件の概念的妥当性や個別領域での経験的妥当性そのものではない。

検証コードは https://github.com/karesansui-u/delta-survival-papers/tree/main/lean において公開されている。検証対象は以下の通りである。 - AxiomsToExp.lean: A1--A3 を含む公理環境のもとでの指数表現の導出。独立積モデル（\(jointSurvival = \prod p_{i}\)）を用いて検証されている。本文3節で述べた望遠鏡積による恒等式（A3 を要しない）は、この検証の特殊ケースとして含まれるが、A1--A2 のみの版は別途形式化されていない - WeakDependence.lean: A3 を弱依存に緩和した場合の ρ-付き指数境界 - RobustSurvival.lean: 弱依存条件のもとでの対応する拡張結果

\begin{itemize}
\tightlist
\item
  LogUniqueness.lean: Paper 1 §3 の対数比の一意性定理（B1--B4 から Cauchy の関数方程式を経て対数関数形 \(f(r) = -k \ln r\) を一意に導出）。基底展開は \(g(t) := f(e^{-t})\) の奇関数拡張による \(\mathbb{R}\) 上の加法写像化を経由し、CauchyExponential.lean の連続加法関数の線形性を再利用している。A3 の確率的独立性には依存せず、AxiomsToExp.lean とは独立のファイルとして閉じている。
\end{itemize}

この LogUniqueness.lean の成立により、A2 の関数形は Lean 上でも定義ではなく定理として特徴づけられ、Paper 2 §3 の恒等式は A3（確率的独立性）を必要とせずに、A1 と対数比の一意性定理のみから形式検証可能になっている。

他方で、形式検証は、個別領域における測度の選択、資源量の具体化、経験的妥当性、あるいは応用の成功そのものまでを保証するものではない。ここで確認されるのは、条件を明示したうえでの導出の論理的一貫性である。

\section{限界}\label{ux9650ux754c}

本稿は、構造持続のすべてを一般的に証明するものではない。また、あらゆる系において A1--A3 がそのまま成立すると主張するものでもない。

本稿が与えるのは、より限定された主張である。すなわち、A1 は制約追加モデルでは命題として従い、より一般には適用条件として置かれ、A2 の関数形の特徴づけ（Paper 1 §3）と累積損失の望遠鏡積とから、指数式は恒等式として成り立ち、段階損失間の依存が累積損失の相対誤差として制御される場合には、指数境界として安定に保たれる、ということである。

したがって、本稿の役割は、応用の全体を与えることではなく、最小形式の背後にある条件つきの数学的骨格を、できる限り薄く与えることにある。

\section{結論}\label{ux7d50ux8ad6}

本稿では、構造持続の指数表現が A1--A2 のもとで恒等式として成り立つことを確認し、段階損失の生成過程に弱い依存がある場合の安定性を述べた。A2 の関数形（対数比）は Paper 1 §3 の対数比の一意性定理によって特徴づけられ、その上で指数表現は累積損失の望遠鏡積として現れる。この二段階の導出はいずれも独立性（A3）を要しない。さらに、A3 を弱依存へ緩和した場合にも、実現累積損失の相対誤差が \(\rho\) で制御される限り、残存可能性は \[
e^{-L_{\text{ref}}(1+\rho)} \le P \le e^{-L_{\text{ref}}(1-\rho)}
\]

の形の指数境界に挟まれる。したがって、完全な独立が成り立たない場合にも、指数型は安定に保たれる。

このことにより、最小形式は単なる見通しのよい表現にとどまらず、条件を明示した骨格として読み直すことができる。本稿の役割は、その骨格を必要最小限のかたちで示すことにある。
\end{document}
